\documentclass[12pt]{article}
\usepackage[utf8]{inputenc}
\usepackage[T1]{fontenc}
\usepackage{hyperref}
\title{Umatilla Language Resources for K-12 Curriculum Development}
\author{Compiled by Claude Sonnet 5 (claude-sonnet-5) for Andr\'as Kornai}
\date{}
\begin{document}
\maketitle

\section{Introduction}
Umatilla (ISO 639-3 uma; Glottocode umat1237) is the southern dialect of Sahaptin (Ichishk\'iin), spoken by the Confederated Tribes of the Umatilla Indian Reservation (CTUIR) in Oregon. Few fluent elder speakers remain.

\section{Primary Educational Resources}
The CTUIR Language Program, established in 1996 in the Department of Education, is reported to offer instruction from kindergarten to adult levels (contrary to the ``not known to be taught'' note in the prompt): \url{https://dictionary.ctuir.org/?p=42734} and \url{https://dictionary.ctuir.org/?p=44495}. Language page: \url{https://ctuir.org/about/language/}.

\section{Linguistic and Reference Materials}
\begin{itemize}
\item Umatilla Language Online Dictionary (by Noel Rude and the CTUIR Language Program): \url{https://dictionary.ctuir.org/uma/ciskiin/}
\item Launch announcement: \url{https://cuj.ctuir.org/ctuir-launches-umatilla-language-online-dictionary/}
\item UBC Relational Lexicography pages on the dictionary: \url{https://wiki.ubc.ca/Documentation:RelLex/Umatilla_Language_Online_Dictionary}, \url{https://wiki.ubc.ca/Documentation:RelLex/Umatilla_Dictionary}, \url{https://knowledgebase.arts.ubc.ca/?p=3375}
\item Wikipedia: \url{https://wikipedia.com/wiki/ISO_639:uma}
\end{itemize}
A print edition of the dictionary is reported as sold by the CTUIR; no separate URL was seen.

\section{Religious and Community Resources}
The CTUIR (\url{https://ctuir.org/about/language/}) is the community authority for the language.

\section{Digital and Audio Resources}
The online dictionary is the principal digital resource; whether it includes audio was not verified.

\section{Curriculum Development Context}
The tribe has an existing language program, a substantial dictionary and a tribal linguist, so curriculum work should be done in partnership with the CTUIR Language Program. Graded textbooks and primers were not found in the results.
\end{document}